\chapter{Implementácia}
\label{ch:implementacia}

\section{Použité technológie}

\begin{table}[H]
  \centering
  \caption{Použité technológie}
  \label{tab:technologie}
  \begin{tabularx}{\textwidth}{@{}l X@{}}
    \toprule
    Oblasť & Technológia \\
    \midrule
    Server & Python 3.12, FastAPI, Pydantic, Uvicorn \\
    Databáza & SQLAlchemy 2, Alembic (migrácie), SQLite / PostgreSQL \\
    Biometria & OpenCV (YuNet \cite{wu2023yunet}, SFace \cite{zhong2021sface}), ONNX Runtime (MiniFASNet \cite{minivision2020silentface}) \\
    Bezpečnosť & PyJWT (prístupové tokeny), cryptography/Fernet (šifrovanie vzorov) \\
    Klient & TypeScript, React 19, React Router, Vite \\
    Kvalita kódu & pytest, Ruff, mypy, ESLint, Vitest, GitHub Actions \\
    Nasadenie & Docker, docker compose, nginx, Render (server), Vercel (klient) \\
    \bottomrule
  \end{tabularx}
\end{table}

Všetky neurónové siete sú predtrénované a voľne dostupné. Sťahuje ich skript
\texttt{scripts/download\_models.py}, ktorý zároveň uloží kontrolné súčty SHA-256.
Spolu majú približne 42~MB a bežia na CPU.

\section{Štruktúra projektu}

\begin{lstlisting}
backend/
  app/
    api/            HTTP rozhranie (routy, závislosti)
    biometrics/     detektor, príznaky, kvalita, živosť, engine
      liveness/     pasívna detekcia, aktívna výzva, póza hlavy
    core/           konfigurácia, JWT, šifrovanie, chyby, rate-limit
    db/             ORM modely, pripojenie, migrácie
    repositories/   prístup k dátam
    services/       registrácia, overenie, identifikácia, výzvy
  evaluation/       metriky ISO/IEC, vyhodnocovacie skripty, grafy
  tests/            jednotkové a end-to-end testy
frontend/src/
  api/  auth/  biometrics/  camera/  components/  pages/
docs/               dokumentácia
report/             táto správa (LaTeX)
\end{lstlisting}

\section{Detekcia tváre a kontrola kvality}

Detektor YuNet sa volá cez triedu \texttt{cv2.FaceDetectorYN}. Pre každú tvár vráti
ohraničujúci obdĺžnik, päť orientačných bodov a skóre. Tváre sa zoradia podľa plochy
a ďalej sa pracuje s~najväčšou. Keďže objekt OpenCV si pamätá veľkosť vstupu, prístup
k~nemu chráni zámok, aby boli súbežné požiadavky bezpečné.

Kontrola kvality (trieda \texttt{QualityAssessor}) je inšpirovaná normou ISO/IEC
29794-5 \cite{iso29794-5} a prehľadom \cite{schlett2022fiqa}, ale je zjednodušená na
rýchle, interpretovateľné kritériá:
\begin{itemize}
  \item v~zábere nesmie byť ďalšia tvár väčšia ako 15~\% plochy hlavnej tváre,
  \item kratšia strana tváre aspoň 90~px,
  \item ostrosť – rozptyl Laplaciánu na výreze tváre zmenšenom na $128\times128$~px aspoň 25,
  \item priemerný jas výrezu v~intervale $\langle 50, 215\rangle$.
\end{itemize}
Výsledkom je zoznam problémov (napr. \texttt{too\_dark}), ktoré klient zobrazí ako
konkrétnu radu. Rovnaká kontrola beží aj ako samostatný endpoint
\texttt{/biometrics/quality}, ktorý klient volá každých 1,2~s počas nastavovania tváre.

\section{Pasívna detekcia živosti}

Použité sú dva modely MiniFASNet (V2 a V1SE) z~projektu Silent-Face-Anti-Spoofing
\cite{minivision2020silentface}, exportované do formátu ONNX. Každý model dostane výrez
tváre zväčšený $2{,}7\times$ resp. $4\times$ oproti ohraničujúcemu obdĺžniku, aby videl
aj okolie tváre (rám displeja, okraj papiera), zmenšený na $80\times80$~px. Vstupom je
obraz BGR s~hodnotami 0--255 vo formáte NCHW, výstupom softmax nad troma triedami,
z~ktorých trieda~1 znamená živú tvár. Výsledné skóre je vážený priemer oboch modelov:
\begin{equation}
  p_{\mathrm{live}} = \frac{\sum_{m} w_m\, p_m(\text{live})}{\sum_m w_m}.
\end{equation}

\begin{lstlisting}[language=Python, caption={Predspracovanie a inferencia jedného modelu (zjednodušené)}]
crop = crop_with_context(image, face.bbox, scale, (80, 80))
tensor = np.transpose(crop.astype(np.float32), (2, 0, 1))[None]
logits = session.run(None, {input_name: tensor})[0]
p_live = softmax(logits)[0, LIVE_CLASS_INDEX]
\end{lstlisting}

\section{Aktívna výzva}

Výzvy vydáva a overuje trieda \texttt{ChallengeService}. Výzva sa uloží pod náhodným
identifikátorom (\texttt{secrets.token\_urlsafe(24)}) s~dobou platnosti 90~s. Pri
overení sa výzva z~úložiska odstráni atomickou operáciou (v~Redise \texttt{GETDEL}),
takže tú istú výzvu nemožno použiť dvakrát. Overuje sa aj účel výzvy -- výzva vydaná na
prihlásenie nemôže byť použitá na registráciu.

Klient každý krok výzvy zobrazí ako pokyn so šípkou a ukazovateľom priebehu. Po
úvodnej pauze 700~ms zachytí päť rovnomerne rozložených snímok počas 2,2~s. Snímky sa
z~videa kopírujú cez \texttt{canvas} bez zrkadlenia, zmenšia na šírku 640~px a odošlú ako
JPEG v~kódovaní base64. Náhľad v~prehliadači je zrkadlený, aby pokyn \enquote{doľava}
zodpovedal smeru, ktorý používateľ vidí.

Vyhodnotenie výzvy na serveri (trieda \texttt{ActiveLivenessVerifier}) vypočíta pre
každú snímku polohu hlavy podľa vzťahov z~kapitoly~\ref{ch:navrh} a pohyb oproti
neutrálnej polohe v~smere požadovanej akcie. Výsledok každého kroku (splnený,
počet použiteľných snímok, maximálny pohyb) sa vracia v~odpovedi, čo zjednodušuje
ladenie prahov.

\section{Registrácia, overenie a identifikácia}

Prípady použitia implementuje trieda \texttt{BiometricAuthService}. Všetky tri majú
rovnaký priebeh: spotrebovanie výzvy, kontrola, že snímky pokrývajú všetky kroky,
dekódovanie obrázkov s~limitom veľkosti, vyhodnotenie relácie a nakoniec akcia
podľa prípadu použitia:
\begin{itemize}
  \item \textbf{registrácia} -- kontrola obsadenosti mena, vyhľadanie tváre v~galérii
        (odmietnutie duplicitnej registrácie), uloženie používateľa a šifrovaného vzoru;
  \item \textbf{verifikácia 1:1} -- kontrola dočasného zablokovania, porovnanie
        s~vzormi používateľa; neexistujúci používateľ vráti rovnakú chybu ako nezhoda;
  \item \textbf{identifikácia 1:N} -- porovnanie so všetkými vzormi aktívnych
        používateľov a výber najlepšej zhody nad prahom;
  \item \textbf{aktualizácia vzoru} -- nový vzor sa uloží, iba ak sa zhoduje
        s~pôvodným, čo bráni prevzatiu účtu po krádeži tokenu.
\end{itemize}

Každý pokus sa zapíše do auditného záznamu (režim, výsledok, dôvod zamietnutia,
skóre zhody a živosti, výsledok výzvy, počet snímok, čas spracovania, IP adresa).
Záznam sa zapisuje v~samostatnej transakcii až po potvrdení alebo zrušení hlavnej
transakcie, takže zostane zachovaný aj pri neúspechu. Z~auditu sa počíta aj
dočasné zablokovanie účtu.

\section{Bezpečnosť a ochrana údajov}

\begin{itemize}
  \item \textbf{Šifrovanie vzorov} -- vektor 128 hodnôt typu \texttt{float32} sa šifruje
        schémou Fernet. Kľúč je v~premenných prostredia, nie v~databáze. Vzor je
        viazaný na verziu modelu, po výmene modelu sa starý vzor nepoužije.
  \item \textbf{Prístupové tokeny} -- JWT podpísané HMAC-SHA256 s~platnosťou 30~min,
        povinné položky \texttt{sub}, \texttt{exp}, \texttt{iat}, \texttt{typ}.
  \item \textbf{Ochrana proti zneužitiu} -- limit požiadaviek na IP adresu, limit veľkosti
        požiadavky (20~MB), snímky (2~MB) a počtu snímok (40); obrázky sa zmenšia na
        najviac 1280~px.
  \item \textbf{Produkčný režim} -- aplikácia odmietne štartovať bez nastavených
        tajomstiev; v~odpovediach sa nevracajú skóre, aby útočník nemohol vstup
        iteratívne optimalizovať; vypnutá je dokumentácia API.
  \item \textbf{Klient} -- prísna politika CSP, zákaz vkladania do rámcov, token v~úložisku
        \texttt{sessionStorage} (zanikne zatvorením karty).
  \item \textbf{GDPR} -- používateľ môže vymazať účet aj vzor; pri registrácii je
        informovaný o~spracovaní biometrického údaja.
\end{itemize}

\section{Klientska aplikácia}

Klient má päť obrazoviek: úvod, registrácia, prihlásenie (prepínač 1:1 / 1:N),
používateľský účet a administrácia. Jadrom je komponent \texttt{BiometricCapture},
ktorý spája tri časti:
\begin{itemize}
  \item \texttt{useCamera} -- spustenie a ukončenie kamery, zrozumiteľné chybové
        hlásenia (zamietnutý prístup, chýbajúca kamera, nezabezpečené pripojenie);
  \item \texttt{useQualityMonitor} -- periodická kontrola kvality a farebné označenie
        oválu (zelený = pripravené);
  \item \texttt{useLivenessCapture} -- získanie výzvy, časovanie krokov, zachytenie snímok
        a odoslanie; podporuje zrušenie.
\end{itemize}
Chyby zo servera sa podľa kódu a dôvodu zamietnutia prekladajú na slovenské hlásenia
s~radami, ako postupovať.

\todo{Sem vložte 2--3 snímky obrazovky (registrácia s~pokynom, úspešné prihlásenie,
zamietnutý útok fotografiou) do priečinka \texttt{figures/} a odkomentujte obrázok nižšie.}
% \begin{figure}[H]
%   \centering
%   \includegraphics[width=0.8\textwidth]{screenshot_enroll.png}
%   \caption{Registrácia – pokyn aktívnej výzvy}
% \end{figure}

\section{Nasadenie}

Server aj klient majú vlastný \texttt{Dockerfile}. Lokálne sa celý systém spustí
príkazom \texttt{docker compose up}: kontajner nginx poskytuje statické súbory klienta
a presmerúva \texttt{/api} na server, server pri štarte spustí databázové migrácie
a Redis zdieľa výzvy medzi procesmi. Verejná inštancia beží na službách Render
(server s~PostgreSQL, definícia v~\texttt{render.yaml}) a Vercel (klient,
\texttt{vercel.json}), pričom Vercel presmerúva \texttt{/api} na server -- prehliadač
tak komunikuje iba s~jednou doménou.

\begin{itemize}
  \item Aplikácia: \url{https://face-auth-antispoofing.vercel.app}
  \item Repozitár: \url{https://github.com/opoiasnik/face-auth-antispoofing}
\end{itemize}
